\section{Local prediction with unknown entropy
dynamics}\label{local_prediction-local-prediction-with-unknown-entropy-dynamics}

Part of the \href{DRAFT.md}{manuscript}. This chapter gives an explicit
common-risk family, a uniform direction-exit proof, and statistical
necessity and sufficiency within that family. Numbering L.1--L.8 is
local to this chapter. The model is scalar, the step size is fixed, and
the calibration and future adaptation samples are independent.

\subsection{8B.1. Observation
experiment}\label{local_prediction-b.1.-observation-experiment}

Write s(u)=(1+exp(-u))\^{}\{-1\}, a(u)=s(u)(1-s(u)), and \[
U(x,z)=-w^2(z-x)a(w(z-x)),\quad
x_{t+1}=x_t+\frac\eta n\sum_{i=1}^n U(x_t,Z_{t,i}),\quad x_0=0.
\] Fix w\textgreater0. Future batches are fresh iid scores from an
unknown environment Q. A predictor observes only M independent
calibration scores, x0, the common model family and its constants, and,
if desired, the entire common classification-risk function. Future SGD
samples, trajectories, root, and environment identity are unavailable.
Any source information is fixed across environments. Randomized
predictors may use independent randomization.

For a positive fixed tolerance epsilon define \[
T_\epsilon=\inf\{t\ge0:R_Q(x_t)>R_Q(0)+\epsilon\}.
\] Below the safe set is exactly {[}-b,b{]}, T\_epsilon is finite almost
surely, and \[
\pi_n(Q)=\Pr_Q(x_{T_\epsilon}>b).
\] The probability integrates over independent future SGD; prediction
loss averages over calibration randomness. Online trajectory prediction
is a different experiment.

Known exact centering admits the asymptotically correct constant
forecast .5. Thus the centered plug-in error arctan(r)/pi in
\href{calibration.md}{calibration} is not a universal lower bound. Here
initialization stays fixed while the unknown environment moves the root.

\subsection{8B.2. Explicit common-risk
family}\label{local_prediction-b.2.-explicit-common-risk-family}

\subsubsection{Proposition L.1. Legal environments and local
geometry}\label{local_prediction-proposition-l.1.-legal-environments-and-local-geometry}

Fix nu\textgreater0 and 0\textless kappa\textless1. Let f0 be the
N(0,nu²) density, and put \[
v_0(x)=E_0 U(x,Z),\quad \lambda=v_0'(0),\quad
\sigma_g^2=E_0 U(0,Z)^2,\quad h(z)=U(0,z)/\sigma_g.
\] Then \[
\lambda=\frac{w^2}{\nu^2}E_0[Z^2a(wZ)]>0,\qquad \sigma_g>0,
\quad E_0h=0,\quad E_0h^2=1,\quad H_*:=\|h\|_\infty<\infty.
\] Choose a sufficiently small d\textgreater0 such that
v0'(x)\textgreater=lambda/2 on K={[}-d,d{]}, and set b=d/2. Choose and
then fix \[
0<\eta<\min\{1/L_U,d/(4B_U)\},\quad B_U=w/4,\quad L_U=w^2/4.
\] There exists bar\_theta\textgreater0 such that for
\textbar theta\textbar\textless=bar\_theta the joint densities \[
f_\theta=f_0(1+\theta h),\qquad g(z)=\kappa\tanh(z)f_0(z),
\qquad p_\theta(z,\pm1)=(f_\theta(z)\pm g(z))/2
\] are nonnegative and normalized, with equal class priors. The score
law Q\_theta has density f\_theta. All environments share the entire
risk function \[
R(x)=\frac12+\int_{-\infty}^x g(z)\,dz,
\quad R(x)-R(0)=\kappa\int_0^x\tanh(z)f_0(z)\,dz.
\] In particular, with the fixed choice \[
\epsilon=\kappa\int_0^b\tanh(z)f_0(z)\,dz>0,
\quad \{x:R(x)\le R(0)+\epsilon\}=[-b,b].
\] The field v\_theta has exactly one root q\_theta in K, with
\textbar q\_theta\textbar\textless=b/2 and
v\_theta'\textgreater=lambda/4 throughout K. Moreover, \[
q_\theta=-\frac{\sigma_g}{\lambda}\theta+O(\theta^2).
\tag{L.1}
\] All constants and the remainder are uniform on a sufficiently small
parameter interval.

\textbf{Proof.} Symmetry makes U(0,z) odd and v0(0)=0. Since U(x,z) is a
translate of U(0,z), Gaussian integration by parts gives
v0'(0)=-E0{[}partial\_z U(0,Z){]}=-E0{[}Z U(0,Z){]}/nu², yielding the
displayed positive formula. The variance is positive because U(0,Z) is
not constant. Boundedness and smoothness follow from Lemma E.1 in
\href{scalar_exit.md}{scalar exit}.

Require bar\_theta H\_*\textless=min(1/2,1-kappa). Then
f\_theta\textgreater=kappa f0\textgreater=\textbar g\textbar, and E0h=0
normalizes f\_theta. Since g is odd, its integral is zero and each class
has mass 1/2. Integrating the two misclassification regions gives the
displayed risk. Its derivative g is negative on the negative half-line
and positive on the positive half-line, so its even safe set is exactly
the stated interval, globally. Thus risk exit equals interval exit path
by path, including discrete overshoots; there is no second safe
component.

Write v\_theta=v0+theta k, k(x)=E0{[}U(x,Z)h(Z){]}. All fixed-order
derivatives of k are bounded, and k(0)=sigma\_g. Shrinking bar\_theta
gives v\_theta'\textgreater=lambda/4 on K and opposite signs at +/-b/2.
The implicit-function theorem gives (L.1), with a bounded second
derivative on the reduced closed parameter interval. Root noise variance
is continuous and stays bounded below by a positive constant after
another reduction. The full one-step overshoot interval {[}-b-eta
B\_U,b+eta B\_U{]} lies in {[}-3d/4,3d/4{]}, inside K. These choices
establish common geometry rather than assuming it. This family is a
general scalar joint model; its label-conditionals are not the original
Gaussian mixture. \(\square\)

\subsection{8B.3. Uniform direction exit for actual entropy
SGD}\label{local_prediction-b.3.-uniform-direction-exit-for-actual-entropy-sgd}

Define \[
\lambda_\theta=v_\theta'(q_\theta),\quad A_\theta=1+\eta\lambda_\theta,
\quad \sigma_\theta^2=\operatorname{Var}_\theta U(q_\theta,Z),
\quad s_\theta^2=\frac{\eta^2\sigma_\theta^2}{A_\theta^2-1}.
\] At theta=0 write A=1+eta lambda, s=s0, and \[
c=\frac{\sigma_g}{\lambda s}=\sqrt{\frac{2+\eta\lambda}{\eta\lambda}},
\qquad e_n=n^{-1/2}+((1+\log n)/n)^{1/3}.
\]

\subsubsection{Lemma L.2. Uniform local direction
law}\label{local_prediction-lemma-l.2.-uniform-local-direction-law}

For every fixed finite A0\textgreater0 there are C and n0 such that \[
\sup_{|u|\le A_0}|\pi_n(Q_{u/\sqrt n})-\Phi(cu)|\le C e_n
\quad(n\ge n_0).
\tag{L.2}
\] This statement concerns fixed eta; its constants are not uniform as
eta tends to zero.

\textbf{Proof.} We supply the common constants and the stopped-seed
argument, so this is not an application of a pointwise exit limit to
changing laws.

Construct normalized C² conjugacies H\_theta with H\_theta(q\_theta)=0,
H\_theta'(q\_theta)=1 and H\_theta(G\_theta(x))=A\_theta H\_theta(x)
whenever both arguments are in K. The uniform inverse construction is as
follows.

Let K={[}-d,d{]} and G\_theta(x)=x+eta v\_theta(x). The signs of
v\_theta at the endpoints imply G\_theta(K) contains K. Thus
F\_theta=G\_theta\^{}\{-1\} is defined on K, maps K into itself, and
fixes q\_theta. Its derivative satisfies

\[
0<F_\theta'(x)\le \gamma:=(1+\eta\lambda/4)^{-1}<1.
\]

In centered coordinates write

\[
\widetilde F_\theta(y)=A_\theta^{-1}y+r_\theta(y),
\qquad r_\theta(0)=r_\theta'(0)=0.
\]

On a common small neighborhood, uniform smoothness gives bounded third
derivatives and a bound beta on the inverse derivative, with A\_max
beta\^{}2\textless1. If y\_j=F\_tilde\_theta\^{}j(y), then, uniformly in
theta,

\[
|y_j|\le C\beta^j|y|,\qquad
|y_j'|\le C\beta^j,\qquad |y_j''|\le C\beta^j.
\]

For the second inequality use the derivative product. The third follows
from

\[
y_{j+1}''=\widetilde F_\theta''(y_j)(y_j')^2
             +\widetilde F_\theta'(y_j)y_j''.
\]

The successive differences in A\_theta\^{}j y\_j equal
A\_theta\^{}\{j+1\} r\_theta(y\_j). Their functions, first derivatives
and second derivatives have summable bounds C(A\_max
beta\textsuperscript{2)}j. This proves local C\^{}2 convergence. To
establish positivity of the limiting derivative, use its derivative
product:

\[
\prod_{k=0}^{j-1} A_\theta\widetilde F_\theta'(y_k).
\]

Each factor is positive and differs from one by
O(\textbar y\_k\textbar). On a sufficiently small common neighborhood
the sum of these deviations is uniformly bounded, so the limiting
product is bounded away from zero. A fixed number of inverse iterations
maps all of K into that neighborhood, since F\_theta contracts by gamma.
Finite composition gives uniform C\^{}2 bounds and a positive derivative
lower bound on K. This also verifies the functional equation wherever
its two arguments are in K.

Thus, for constants independent of theta, \[
\|H_\theta\|_{C^2(K)}\le C,\quad \inf_K H_\theta'\ge m>0,
\quad 1<A_{\min}\le A_\theta\le A_{\max},\quad
0<s_{\min}\le s_\theta\le s_{\max}.
\tag{L.3}
\] The inverse derivative bounds use the uniform smoothness of U and the
positive derivative of G\_theta. The root-to-boundary distances are at
least b/2. These facts also give uniform bounds above and below for the
absolute conjugate boundary values.

Let tau be exit from {[}-b,b{]} and xi\_j(x)=n\^{}\{-1\}sum\_i
U(x,Z\_\{j,i\})-v\_theta(x). Each theta has finite tau almost surely:
choose two fixed bounded score intervals respectively left and right of
K. Their masses are uniformly positive since f\_theta\textgreater=f0/2.
Repeated batches entirely in the left interval drive any state at least
q\_theta rightward by a positive minimum speed until exit; the right
interval does the converse. Compactness supplies a fixed finite block
length N. The conditional chance of an appropriate exiting block is at
least p\^{}\{nN\}\textgreater0; geometric blocking proves finiteness.
This is not a bound on the typical exit time.

For \textbar theta\textbar\textless=A0/sqrt(n), (L.1) implies
\textbar q\_theta\textbar\textless=C/sqrt(n) and mu=sqrt(n)H\_theta(0)
is bounded. Conditional batch moments give \[
E[\xi_j(x)^2\mid\mathcal F_j]\le C/n,\quad
E[\xi_j(x)^4\mid\mathcal F_j]\le C/n^2,\quad
E[(\xi_j(x)-\xi_j(q_\theta))^2\mid\mathcal F_j]
\le C|x-q_\theta|^2/n.
\] All constants are distribution independent, by bounded U and
partial\_x U. Taylor expansion on the established overshoot domain
gives, for j\textless tau, \[
H_\theta(x_{j+1})=A_\theta H_\theta(x_j)
 +\eta H_\theta'(G_\theta(x_j))\xi_j(x_j)+r_j,
\quad |r_j|\le C\xi_j(x_j)^2.
\] Set Y\_t=sqrt(n) A\_theta\^{}\{-(t wedge tau)\}H\_theta(x\_\{t wedge
tau\}).

Let F\_j denote the history before batch j and write
I\_j=1\_\{j\textless tau\}. With

\[
D_j=\sqrt n\,\eta A_\theta^{-j-1} I_j
 H_\theta'(G_\theta(x_j))\xi_j(x_j),
\]

and

\[
E_j=\sqrt n A_\theta^{-j-1}I_j r_j,
\]

the stopped recursion gives

\[
Y_t=\mu+\sum_{j<t}D_j+\sum_{j<t}E_j.
\]

The indicator and the coefficient are F\_j-measurable. Thus D\_j are
martingale differences and

\[
\sum_j\mathbb E D_j^2\le C\sum_j A_\theta^{-2j}\le C.
\]

The conditional fourth-moment bound gives

\[
\|E_j\|_2\le C n^{-1/2}A_\theta^{-j},
\qquad \left\|\sum_{j<t}E_j\right\|_2\le Cn^{-1/2}.
\]

Therefore sup\_t E\textbar Y\_t\textbar\^{}2\textless=C without assuming
a logarithmic exit-time estimate. Finite exit and Fatou then imply
E\textbar Y\_tau\textbar\^{}2\textless=C. At exit, the uniform
root-to-boundary separation and positive derivative of H imply
\textbar H\_theta(x\_tau)\textbar\textgreater=b\_H\textgreater0.
Consequently

\[
\mathbb E A_\theta^{-2\tau}\le (nb_H^2)^{-1}\mathbb E|Y_\tau|^2\le C/n.
\]

There is no circular inference from a desired exit approximation to its
own moment estimate.

Before exit, H\_theta'\textgreater=m and bounded K also imply \[
E[\mathbf1_{j<\tau}|x_j-q_\theta|^2]\le C\min(A_\theta^{2j}/n,1).
\tag{L.4}
\]

On the same infinite sample array introduce the frozen seed \[
W=\mu+\eta\sum_{j\ge0}A_\theta^{-j-1}\sqrt n\,\xi_j(q_\theta).
\] The difference H\_theta'(G\_theta(x\_j))xi\_j(x\_j)-xi\_j(q\_theta)
has conditional mean zero and variance at most
C\textbar x\_j-q\_theta\textbar²/n, using H\_theta'(q\_theta)=1 and
(L.3). Consequently its stopped weighted sum has second moment at most
\[
C\sum_{j\ge0}A_\theta^{-2j}\min(A_\theta^{2j}/n,1)
\le C(1+\log n)/n.
\] Use the convention that batch j moves x\_j to x\_\{j+1\}. At tau, all
batches with indices j\textgreater=tau remain fresh. Conditional on
F\_tau, the frozen tail has mean zero and variance

\[
\eta^2\sigma_\theta^2\sum_{j\ge\tau}A_\theta^{-2j-2}
=s_\theta^2 A_\theta^{-2\tau}.
\]

This proves the O(1/n) tail second moment in (L.5) below. It is a
stopped-tail independence argument, not a claim that the stopped
trajectory is independent of the full frozen seed.

Together, \[
E|Y_\tau-W|^2\le C(1+\log n)/n.
\tag{L.5}
\] Grouping W-mu by sample column gives n\^{}\{-1/2\}sum\_i V\_i with
iid centered V\_i=eta sum\_\{j\textgreater=0\}
A\_theta\^{}\{-j-1\}U(q\_theta,Z\_\{j,i\}). These variables have uniform
absolute bounds eta B\_U/(A\_min-1), and variance s\_theta² bounded
below. Berry--Esseen gives a uniform C/sqrt(n) Kolmogorov error relative
to N(mu,s\_theta²). In particular
P(\textbar W\textbar\textless=delta)\textless=C delta+C/sqrt(n).

The sign of Y\_tau is exactly the exit direction. Hence (L.5),
Chebyshev, and the preceding anti-concentration bound give \[
|\pi_n(Q_\theta)-\Phi(\mu/s_\theta)|
\le C\{n^{-1/2}+\delta+(1+\log n)/(n\delta^2)\}\le C e_n,
\] choosing delta=((1+log n)/n)\^{}\{1/3\}. This explicitly controls
seeds near zero, rather than identifying their signs without an error
bound. Finally normalization and the uniform C² bound imply
H\_theta(0)=-q\_theta+O(q\_theta²); smoothness gives
s\_theta=s+O(\textbar theta\textbar). Uniformly on
\textbar u\textbar\textless=A0, \[
\mu/s_\theta=cu+O(n^{-1/2}).
\] The Lipschitz property of Phi proves (L.2). \(\square\)

\subsection{8B.4. Statistical threshold within the same
family}\label{local_prediction-b.4.-statistical-threshold-within-the-same-family}

Let \[
\mathfrak R_{n,M}(A_0)=\inf_{\widehat\pi}
\sup_{|u|\le A_0}E_{u/\sqrt n}
|\widehat\pi-\pi_n(Q_{u/\sqrt n})|,
\] where A0\textgreater0 and predictors have the information in §8B.1.
The family and its constants are known; u is unknown.

\subsubsection{Theorem L.3. Necessity with explicit
constants}\label{local_prediction-theorem-l.3.-necessity-with-explicit-constants}

If eventually M/n\textless=C0\textless infinity, choose \[
0<a\le\min\{A_0,1/(4\sqrt{\max(C_0,1)})\}.
\] Then \[
\liminf_n\mathfrak R_{n,M}(A_0)\ge [2\Phi(ca)-1]/4>0.
\tag{L.6}
\]

\textbf{Proof.} At theta=+/-a/sqrt(n), positivity gives \[
D_{KL}(f_\theta\|f_{-\theta})\le\chi^2(f_\theta,f_{-\theta})
=4\theta^2 E_0\frac{h^2}{1-\theta h}\le8\theta^2.
\] Product KL is at most 8a²M/n.~Pinsker gives total variation at most
2a sqrt(M/n)\textless=1/2. Lemma L.2 implies the target separation
Delta\_n tends to Delta=2Phi(ca)-1\textgreater0. Let p\_+,p\_- be the
calibration densities and t\_+,t\_- their true exit probabilities. For
any forecast T, the triangle inequality on their common density overlap
gives \[
E_+|T-t_+|+E_-|T-t_-|
\ge\int\min(p_+,p_-)(|T-t_+|+|T-t_-|)
\ge\Delta_n(1-\operatorname{TV}).
\] The maximum risk is at least half this sum, hence at least Delta\_n/4
when TV\textless=1/2. This proves (L.6). Common risk and fixed source
information have identical values and cannot change this comparison.
Independent randomization cannot improve the testing bound. \(\square\)

\subsubsection{Theorem L.4. Sufficiency and
equivalence}\label{local_prediction-theorem-l.4.-sufficiency-and-equivalence}

The predictor \[
\widehat\theta=M^{-1}\sum_{j=1}^M h(\widetilde Z_j),\qquad
\widehat\pi=\Phi(c\sqrt n\,\widehat\theta)
\] satisfies, for M\textgreater=1 and large n, \[
\sup_{|u|\le A_0} E_{u/\sqrt n}|\widehat\pi-\pi_n(Q_{u/\sqrt n})|
\le \frac{c}{\sqrt{2\pi}}\sqrt{n/M}+Ce_n.
\tag{L.7}
\] Thus uniform consistency within this family is possible if and only
if M/n tends to infinity.

\textbf{Proof.} E\_theta h=theta exactly, and Var\_theta
h=1-theta²\textless=1 by oddness of h (see L.5). Hence
E\textbar hat\_theta-theta\textbar\textless=1/sqrt(M). Apply the global
Lipschitz constant 1/sqrt(2pi) of Phi and (L.2). No clipping or root
oracle is used. Conversely, if M/n does not diverge, some infinite
subsequence has bounded M/n; Theorem L.3 on that subsequence excludes
risk tending to zero. \(\square\)

These consistency statements are specific to the known family and fixed
eta. The following refinements distinguish minimax conclusions from the
risk of an individual predictor.

Retain the notation of L.1--L.4. Let A0\textgreater0 be fixed, let

\[
r_n=\sup_{|u|\le A_0}|\pi_n(Q_{u/\sqrt n})-\Phi(cu)|\le C e_n,
\qquad B_0=\Phi(cA_0)-\tfrac12.
\]

The quantity r\_n is a deterministic approximation error, not an
observable calibrated confidence bound.

\subsubsection{Proposition L.5. Exact symmetry and calibration
moments}\label{local_prediction-proposition-l.5.-exact-symmetry-and-calibration-moments}

For every legal theta,

\[
\mathbb E_\theta h=\theta,\qquad
\mathbb E_\theta h^2=1,\qquad
\operatorname{Var}_\theta h=1-\theta^2.
\tag{L.8}
\]

For every finite batch size n,

\[
\pi_n(Q_{-\theta})=1-\pi_n(Q_\theta),
\qquad \pi_n(Q_0)=\tfrac12.
\tag{L.9}
\]

In particular the moment predictor in L.4 satisfies the sharper bound

\[
\sup_{|u|\le A_0}\mathbb E_{u/\sqrt n}
 |\widehat\pi-\pi_n(Q_{u/\sqrt n})|
\le\frac{c}{\sqrt{2\pi}}\sqrt{n/M}+r_n.
\tag{L.10}
\]

\textbf{Proof.} The function h is odd and f0 is even. Therefore E0
h\^{}3=0, and E\_theta h\^{}2=E0 h\^{}2+theta E0 h\^{}3=1. The first
identity was already established in L.4, giving the variance formula. In
particular
E\textbar hat\_theta-theta\textbar\textless=sqrt((1-theta\^{}2)/M)\textless=1/sqrt(M);
apply the Lipschitz constant of Phi.

For symmetry, f\_\{-theta\}(z)=f\_theta(-z) and U(-x,-z)=-U(x,z).
Reflect every score in a future sample array. Starting at zero, the
entire reflected path is the negative of the original path. The
symmetric safe interval has the same exit time and the opposite exit
side. Since exit is almost surely finite, the two directional
probabilities sum to one. This proves (L.9). The exact equality at
theta=0 belongs to this symmetric constructed family; it should not be
transferred to asymmetric visual safe sets. \(\square\)

\subsubsection{Proposition L.6. Stronger absolute-loss lower
bound}\label{local_prediction-proposition-l.6.-stronger-absolute-loss-lower-bound}

For all sufficiently large n and any 0\textless a\textless=A0, put

\[
\Delta_0(a)=2\Phi(ca)-1,\qquad \rho_n=M/n.
\]

Then

\[
\mathfrak R_{n,M}(A_0)
\ge \frac{[\Delta_0(a)-2r_n]_+}{2}
       [1-2a\sqrt{\rho_n}]_+.
\tag{L.11}
\]

Consequently, under the eventual bound and choice of a in L.3, this
recovers

\[
\liminf_n\mathfrak R_{n,M}(A_0)
\ge [2\Phi(ca)-1]/4.
\tag{L.12}
\]

\textbf{Proof.} Let P\_+,P\_- be the two calibration laws and t\_+,t\_-
their true SGD direction probabilities. Their densities can be taken
relative to a common measure. For any forecast T,

\[
\begin{aligned}
\mathbb E_+|T-t_+|+\mathbb E_-|T-t_-|
&\ge\int\min(p_+,p_-)\bigl(|T-t_+|+|T-t_-|\bigr)\\
&\ge |t_+-t_-|\,[1-\operatorname{TV}(P_+,P_-)].
\end{aligned}
\]

Independent randomization can be included in the common measure and
leaves TV unchanged. The maximum of the two risks is at least half their
sum. L.2 implies \textbar t\_+-t\_-\textbar\textgreater=Delta0(a)-2r\_n,
and the KL/Pinsker calculation in L.3 gives TV\textless=2a sqrt(rho\_n).
This proves (L.11). Under TV\textless=1/2 the bound is Delta\_n/4, as
used in L.3. A nearest-target testing reduction would lose a factor of
two. \(\square\)

\textbf{Rate consequence with an explicit remainder.} Set
a\_\emph{=min(A0,1/4) and t\_\{n,M\}=min(1,sqrt(n/M)) for
M\textgreater=1. Apply (L.11) with a=a\_} t\_\{n,M\}. Since

\[
\Delta_0(a)\ge 2c\varphi(cA_0)a,
\]

we obtain

\[
\mathfrak R_{n,M}(A_0)
\ge \frac{c\varphi(cA_0)a_*}{2}\,t_{n,M}-\frac{r_n}{2}.
\tag{L.13}
\]

Together with (L.10) and the constant predictor,

\[
\mathfrak R_{n,M}(A_0)
\le\min\left\{B_0+r_n,
\frac{c}{\sqrt{2\pi}}\sqrt{n/M}+r_n\right\}.
\tag{L.14}
\]

Thus matching order follows in sequences where e\_n=o(t\_\{n,M\}); this
is not a global matching-rate theorem for arbitrarily large M. Once the
estimation scale is below the exit-approximation remainder, these bounds
no longer determine the rate.

\subsubsection{Theorem L.7. Exact minimax limit when M/n tends to
zero}\label{local_prediction-theorem-l.7.-exact-minimax-limit-when-mn-tends-to-zero}

If M/n-\textgreater0, including bounded M, then

\[
\boxed{\lim_{n\to\infty}\mathfrak R_{n,M}(A_0)
=\Phi(cA_0)-\tfrac12.}
\tag{L.15}
\]

The constant forecast 1/2 is asymptotically minimax in this regime.

\textbf{Proof.} For the upper bound, L.2 and monotonicity of Phi give

\[
\sup_{|u|\le A_0}|\tfrac12-\pi_n(Q_{u/\sqrt n})|
\le B_0+r_n.
\]

For the lower bound take a=A0 in (L.11). The calibration TV bound tends
to zero and the target separation tends to 2B0, so liminf
risk\textgreater=B0. These two bounds agree. If M=0 is allowed, use the
identical no-observation experiment, with TV exactly zero. \(\square\)

This is an exact asymptotic constant only in the regime
M/n-\textgreater0. It is not a solution of the positive finite-ratio
minimax problem. The result strengthens the interpretation of the sample
threshold: a best constant predictor may be optimal when information
vanishes, while its nonzero worst-case error still prevents uniform
consistency.

\subsubsection{Theorem L.8. Uniform risk limit of the raw moment
predictor}\label{local_prediction-theorem-l.8.-uniform-risk-limit-of-the-raw-moment-predictor}

Suppose M/n-\textgreater rho in (0,infinity). Let

\[
\widehat\pi^{\rm mom}=\Phi(c\sqrt n\,\widehat\theta),
\qquad r=c/\sqrt\rho.
\]

Define

\[
D(t,r)=\mathbb E|\Phi(t+rZ)-\Phi(t)|,
\qquad Z\sim N(0,1).
\]

Then

\[
\sup_{|u|\le A_0}
\left|\mathbb E_{u/\sqrt n}
 |\widehat\pi^{\rm mom}-\pi_n(Q_{u/\sqrt n})|
-D(cu,r)\right|\longrightarrow0.
\tag{L.16}
\]

Moreover, D is even in t, is strictly decreasing for t\textgreater0 when
r\textgreater0, and satisfies

\[
\boxed{\lim_n\sup_{|u|\le A_0}
\mathbb E_{u/\sqrt n}
 |\widehat\pi^{\rm mom}-\pi_n(Q_{u/\sqrt n})|
=\frac{\arctan(c/\sqrt\rho)}{\pi}.}
\tag{L.17}
\]

\textbf{Proof.} By (L.8), the standardized calibration sum

\[
Z_{n,u}=\frac{\sqrt M(\widehat\theta-u/\sqrt n)}
 {\sqrt{1-u^2/n}}
\]

has a uniform Berry--Esseen bound C/sqrt(M) on the compact u-window: h
is bounded and its variance is uniformly bounded away from zero for
large n.~Hence

\[
c\sqrt n\widehat\theta
=cu+c\sqrt{n/M}\sqrt{1-u^2/n}\,Z_{n,u}.
\]

The coefficient of Z\_\{n,u\} tends uniformly to r. This implies uniform
convergence of the bounded loss expectations to D(cu,r). One direct
justification is by subsequences: any alleged nonuniformity supplies
u\_n in the compact window; extract u\_n-\textgreater u\_*, apply the
triangular CLT and bounded continuous mapping along that subsequence,
and use continuity of D. Replacing Phi(cu) by the true pi\_n changes
expected absolute loss by at most r\_n.~This proves (L.16).

Pair positive and negative Z to obtain

\[
D(t,r)=\int_0^\infty
 [\Phi(t+rz)-\Phi(t-rz)]\varphi(z)\,dz.
\tag{L.18}
\]

Symmetry makes this even in t. For t\textgreater0 and z\textgreater0,
\textbar t+rz\textbar\textgreater\textbar t-rz\textbar, so
differentiation under the integral gives a strictly negative derivative
in t. Therefore the maximum over a window containing zero is D(0,r).

Finally,

\[
\frac{\partial}{\partial r}D(0,r)
=2\int_0^\infty z\varphi(rz)\varphi(z)\,dz
=\frac1{\pi(1+r^2)}.
\]

Since D(0,0)=0, integration yields D(0,r)=arctan(r)/pi. Uniform
convergence in (L.16) allows passage to the supremum, proving (L.17). \(\square\)

The arctan expression remains a result for this specified predictor. It
is now its worst-case local-window risk limit, rather than a universal
lower bound. In this known family, the constant predictor has worst-case
limiting risk B0; it can outperform the raw moment predictor for some
combinations of window width and calibration ratio.

\subsubsection{A window-aware
alternative}\label{local_prediction-a-window-aware-alternative}

When the parameter window is declared known, define

\[
\widehat u^{\rm clip}=\operatorname{proj}_{[-A_0,A_0]}(\sqrt n\widehat\theta),
\qquad \widehat\pi^{\rm clip}=\Phi(c\widehat u^{\rm clip}).
\]

For each u in the window, projection onto the corresponding probability
interval gives the pointwise inequality

\[
|\widehat\pi^{\rm clip}-\Phi(cu)|
\le |\widehat\pi^{\rm mom}-\Phi(cu)|.
\]

For loss against the exact SGD probability, its expected risk is at most
the raw predictor's risk plus 2r\_n.~It has the same sufficiency
threshold, and is a useful additional comparator. This observation does
not establish its minimax optimality.

\subsection{8B.5. Empirical interpretation and contribution
boundary}\label{local_prediction-b.5.-empirical-interpretation-and-contribution-boundary}

Use the explicit family first. Keep x0=0, w, nu, eta, kappa, b and the
tolerance fixed. Check the common compact geometry before selecting
numerical constants; the phrase ``sufficiently small d'' must be
instantiated for the chosen run.

For an interpretable horizontal axis, choose a grid of seed coordinates
t in {[}-K,K{]} and generate theta=t/(c sqrt(n)). Here the theoretical
u-window is {[}-K/c,K/c{]}. The generator knows theta; the predictor
receives only its calibration scores and common constants. Do not supply
t, u, the true root, or the exact target pool to the predictor.

Exact sampling is possible by proposing Z\textasciitilde f0 and
accepting with probability (1+theta
h(Z))/(1+\textbar theta\textbar H\_*). The proposal procedure is
internal to generation and must not become side information for
prediction.

Use four predictors: constant .5; raw moment; window-clipped moment; and
empirical-root with estimated scale. Define a fallback for every failed
empirical root or scale fit, retaining all failures in the primary risk
denominator. The moment predictor has no root-solving failure.

Choose calibration budgets through the relative uncertainty r=c
sqrt(n/M). For example r in \{2,1,.5,.25\} corresponds to M
approximately c\^{}2 n/r\^{}2. Report both M/n and r. Because c\^{}2 can
be large at a small fixed step size, modest nominal ratios M/n may all
lie in a poor practical regime despite the asymptotic threshold.

Use a separate sequence with M/n increasing to test consistency, and a
sequence with M/n decreasing to test the approach to the
constant-predictor benchmark. A fixed-r sequence tests L.8; it does not
test consistency.

Estimate each true pi\_n from independent future SGD paths. Record
complete paths or at least reproducible random-stream state and exit
records so capped paths can be continued without changing observed
prefixes. With N paths, N\_R right exits and N\_C censored paths, the
eventual right-exit fraction in that sample lies between N\_R/N and
(N\_R+N\_C)/N. Report this censoring bracket separately from Monte Carlo
uncertainty. If an interval for the population probability is needed,
use simultaneous binomial bounds for the fixed-cap events ``right exit
by cap'' and ``not left exit by cap''; the former lower bound and latter
upper bound bracket pi\_n with the declared confidence level. Do not
treat nonexit by the cap as a third permanent outcome.

Use the same independent reference estimate for all predictors in a
cell. If its error is at most delta, every estimated absolute
probability loss differs from its true-target loss by at most delta, by
the reverse triangle inequality. Thus reference Monte Carlo uncertainty
must be smaller than the scale of predictor differences being
interpreted. Conditional on a fixed reference sample, reuse across
calibration repetitions is not additional independent evidence about
pi\_n.

Report the cellwise mean absolute error, calibration-repetition
uncertainty, grid maximum error, root failure frequency, root error, and
censoring. At finite n the moment predictor need not achieve its maximum
exactly at u=0; L.17 is an asymptotic assertion. A finite grid does not
certify uniformity over the continuum.

No sequential stopping rule or gain-selection algorithm is required for
this experiment.

\subsection{8B.6. Contribution and open
questions}\label{local_prediction-b.6.-contribution-and-open-questions}

We study the independent information required to predict the direction
of a future risk exit under entropy adaptation. We construct a scalar
family with fixed initialization and an identical complete
classification-risk function across environments, while the unknown
unlabeled marginal moves an unstable entropy root on the mini-batch
noise scale. A uniform approximation for the actual nonlinear SGD
direction connects this dynamical sensitivity to a statistical
estimation problem. Within this known family and at fixed positive step
size, uniform consistency is possible exactly when the independent
calibration size grows faster than the SGD batch size. When their ratio
tends to zero, the exact limiting minimax risk is attained by the
constant forecast one half. At positive finite ratios, an explicit
moment predictor instead has a nonvanishing worst-case risk limit
governed by an arctangent law; this predictor-specific limit is not a
minimax lower bound.

The finite-ratio problem is continued in §8B.7 below. Global sharp rates
below the exit-approximation remainder, robustness to a misspecified
tilt direction, broader unknown score families, joint vanishing-step
limits and neural-network adaptation remain separate questions.
Literature-priority assessment remains necessary.

\subsection{8B.7. Finite-information minimax
precision}\label{local_prediction-b.7.-finite-information-minimax-precision}

Read the \href{gaussian_decision.md}{Gaussian decision extension} next
as a continuation of this manuscript. It retains the full proofs of
Theorem L.9 and Corollary L.10, with its own G-equation numbering, and
the numerical analysis in its original location. The
\href{../experiments/results/gaussian_decision/README.md}{reproduction
entry} documents the scripts and evidence boundaries.

Theorem L.9 constructs parameter-independent Markov kernels in both
directions and proves asymptotic equivalence, uniformly over the fixed
local window, between the full calibration experiment and
Y=u+Z/sqrt(rho) when M/n tends to rho in (0,infinity). Its
likelihood-ratio L1 control is stronger than the moment-statistic CLT in
L.8. Corollary L.10 transfers bounded absolute loss and the uniform exit
approximation to obtain \[
\mathfrak R_{n,M}(A_0)\longrightarrow
V_G(cA_0,c/\sqrt\rho),\qquad
V_G(K,r)=\inf_d\sup_{|t|\le K} E_t|d(t+rZ)-\Phi(t)|.
\] The extension gives finite-decimal prior lower bounds and
continuous-window rule upper bounds for six Gaussian problems, with
directed probability calculations and crossing brackets. For K=r=1 the
outward six-decimal bracket is {[}0.164401,0.164587{]}, compared with
the raw moment risk .25. The certified candidate rule therefore improves
Gaussian worst-case MAE by at least 34.16\%; this is not pointwise
dominance. It does not require the searched prior to be exactly least
favorable.

The independent entropy-SGD pilot contains ten cells, 10,240 reference
paths and no censored paths. The largest empirical five-point MAEs for
raw moment versus optimized rules are .253404 versus .167207 at n=128
and .267974 versus .178871 at n=512. These empirical maxima support
finite-batch relevance but are neither continuous-window bounds nor
validation of the narrow Gaussian minimax brackets.
Reference-probability uncertainty is additional to conditional paired
calibration standard errors. The full tables and their interpretation
remain in the
\href{../experiments/results/gaussian_decision/analysis.md\#7-results}{original
analysis}.

The minimax convergence is asymptotic: the proof does not numerically
calibrate the experiment deficiency or the SGD exit remainder. The
pilot's geometry quadrature is diagnostic, separate from the directed
Gaussian certificate. Implementation details, prior/rule specifications,
original outputs and raw pilot arrays stay in the original directory.

\subsection{8B.8. Unknown local directions and fixed-design
confirmation}\label{local_prediction-b.8.-unknown-local-directions-and-fixed-design-confirmation}

Continue with \href{unknown_direction.md}{N.1--N.4 and their complete
proofs}. Around the same known baseline, let f\_\{n,k\}=f0(1+k/sqrt(n)),
E0 k=0, \textbar\textbar k\textbar\textbar\_infinity\textless=B\_loc,
and \textbar E0 psi k\textbar\textless=K, with psi=U(0,.)/(lambda s) and
E0 psi²=c². Provided B\_loc\textgreater=K
\textbar\textbar psi\textbar\textbar\_infinity/c², the class contains
k\_t=t psi/c² for every \textbar t\textbar\textless=K. Its finite-ratio
minimax value is the same V\_G(K,c/sqrt(rho)), and its
uniform-consistency threshold remains M/n tending to infinity.

The lower bound restricts to the contained least-favorable subfamily.
The upper bound applies Gaussian rules to the direction-free statistic
sqrt(n) mean(psi). A uniform exit proof and uniform scalar CLT suffice
for this upper bound; no equivalence of the entire infinite-dimensional
experiment is asserted. The sensitivity bound
\textless psi,h\textgreater²/E0 h²\textless=c² identifies the
entropy-gradient direction as least favorable for a fixed seed window.
Known baseline constants, a fixed bounded envelope and n\^{}(-1/2)
locality are essential. The nuisance projection formulas and their
admissibility qualifications remain in N.4.

The
\href{../experiments/results/precision_confirmation/analysis.md}{independent
confirmation analysis} supersedes the pilot as the precision evidence:
it uses the frozen Gaussian rule and ten fixed cells, with simultaneous
intervals that account for both reference trajectories and calibration
repetitions. It tests the original tilt family, not the new
unknown-direction or nuisance class. Its full tables and raw data remain
at their original location. The \href{RELATED_WORK.md}{contribution
comparison} separates these results from prior entropy-collapse,
monitoring and unlabeled action-selection results.
